\subsection{Every fiber of the cuspidally central cubic quotient}

\begin{statement}{E.36}{Proved}{A class-two quotient retaining all ten abelian coordinates, with every fiber evaluated}
Let $E=X_3^+$, let $D$ be its nine coordinate cusps, and put
$O=[1:-1:0]$. Use the tangential base point at $O$ with
parameter $s=z/x$ and tangent one, both on $U_3=E\setminus D$
and on $E\setminus\{O\}$.
Write
\[
 N=\pi_1(U_{3,\overline{\mathbb Q}})^{(3)},\qquad
 V=N^{\rm ab}=T_3J_{E,D},\qquad T=T_3E,
 \qquad I=\ker(N\longrightarrow T).
\]
Closures are understood in commutator subgroups. Set
\[
 N^{\rm cc}=N/[N,I],\qquad
 C=\pi_1((E\setminus\{O\})_{\overline{\mathbb Q}})^{(3)}
                   /\gamma_3 .
\]
The maps induced by the open immersion
$U_3\hookrightarrow E\setminus\{O\}$ and by abelianization
give an actual, based, $G_{\mathbb Q}$-equivariant isomorphism
\[
 N^{\rm cc}\ \xrightarrow{\ \sim\ }\ C\times_T V,
       \qquad n\longmapsto(q_O(n),n^{\rm ab}).
 \tag{E36.1}
\]
In particular the quotient retains all ten coordinates of
$V$ and its nonzero elliptic commutator. The exact quotient
from the full intrinsic class-two group is
\[
 1\longrightarrow
 \ker\!\left(\bigwedge\nolimits^2 V
           \longrightarrow\bigwedge\nolimits^2 T\right)
 \longrightarrow N/\gamma_3N
 \longrightarrow N^{\rm cc}\longrightarrow1 .
 \tag{E36.2}
\]
The displayed kernel is a free $\mathbb Z_3$-module of rank
$44$, with the Galois action induced on the exterior square;
the remaining central commutator module is
$\bigwedge^2T\simeq\mathbb Z_3(1)$.

Every geometric-abelian section class
$\alpha\in H^1(\mathbb Q,V)$ lifts to a section of
$N^{\rm cc}\rtimes G_{\mathbb Q}$.
For every such $\alpha$, the complete fiber of geometric
conjugacy classes is a torsor under
\[
 \mathcal R_{\mathbb P}(\mathbb Z_3)
 =\left\{(a_p)_{p\ {
m rational\ prime}}\in
                \prod_p\mathbb Z_3:
       \substack{\text{for every }m,\ a_p\in3^m\mathbb Z_3\\
                         \text{ for all but finitely many }p}\right\}.
 \tag{E36.3}
\]
The parameter $a=(a_p)$ acts by the central Kummer cocycle
of the compatible classes $\prod_p p^{a_p}$ modulo
$\mathbb Q^{\times3^m}$. These finite products modulo each
$3^m$ define the indicated element of the completion for
multiplication by powers of three.
Thus this is a computation of every fiber over the entire
rank-ten abelian section group evaluated in E.24, including
its ramified classes. The classes obstructed in E.31 also
lift to (E36.1), while their further lifting through (E36.2)
is obstructed by the explicit quotient of E.31.

\emph{Manuscript: \S51.16, folios 733v--736v, scans 172--178,
(371)--(379). The quotient maps (E36.1)--(E36.2) specify
exactly the intrinsic cubic quotient evaluated here.}
\end{statement}
\begin{proof}
First compute the geometric maps integrally. An oriented
presentation of the punctured complex elliptic curve is
\[
 N=\langle a,b,c_O,c_1,\ldots,c_8\mid
              [a,b]c_Oc_1\cdots c_8=1\rangle_{\rm pro\text{-}3}.
 \tag{E36.4}
\]
Here $a,b$ map to a basis of $T$; the $c_d$ are the
peripheral loops. Eliminating $c_O$ gives the free pro-three
group of rank ten. The kernel $I$ is normally generated by
these peripheral loops. Passing to $N/[N,I]$ makes all
$c_d$ central and leaves only the indicated relation.
Consequently the geometric cuspidal kernel is
\[
 I/[N,I]\simeq\mathbb Z_3(1)[D],
 \qquad [a,b]=-\sum_{d\in D}c_d,
 \tag{E36.5}
\]
with its actual Galois action: the cusps are permuted and
each inertia generator has the cyclotomic twist. This action
can also be obtained from the torus
$\operatorname{Res}_{D/\mathbb Q}\mathbb G_m$.
The central quotient is the one defined in
Borne--Emsalem--Stix, \S3.2; their Proposition 3.5 identifies
it with the fundamental group of their torus torsor by
the canonical map from the open curve
\cite{E_BorneEmsalemStix}.

For an entirely explicit check of (E36.5), take the
two-step Lie module with basis
$a,b,c_1,\ldots,c_8,k$ and only bracket $[a,b]=k$.
The law $u*v=u+v+[u,v]/2$ is integral over $\mathbb Z_3$.
Sending $c_O$ to $-k-\sum_{i=1}^8c_i$ satisfies (E36.4).
Conversely its generators and brackets give precisely the
relations of $N/[N,I]$. The two maps are inverse, so there
are no further relations or torsion in the nine-dimensional
central lattice of (E36.5).

The map $q_O$ fills the other eight punctures. On (E36.5)
it is projection to the $O$-coordinate. Abelianization
quotients (E36.5) by the diagonal element
$\delta=\sum_{d\in D}c_d$. The maps on the two geometric
cuspidal kernels are therefore
\[
 \mathbb Z_3(1)[D]\longrightarrow
 \mathbb Z_3(1)\oplus
       \bigl(\mathbb Z_3(1)[D]/\mathbb Z_3(1)\delta\bigr),
 \quad t\longmapsto(t_O,\overline t).
 \tag{E36.6}
\]
This is an isomorphism: in coordinates it sends
$(t_O,t_1,\ldots,t_8)$ to
$(t_O,t_1-t_O,\ldots,t_8-t_O)$, with inverse
$(u,v_1,\ldots,v_8)\mapsto(u,v_1+u,\ldots,v_8+u)$.
It is $G_{\mathbb Q}$-equivariant because $O$ is rational
and the diagonal element is invariant apart from its
cyclotomic twist. Both sides of (E36.1) map onto $T$,
and (E36.6) identifies their kernels. This proves (E36.1).
The isomorphism is induced by the stated morphisms of
curves and abelianization, with the same tangent throughout.
It therefore identifies the actual arithmetic extensions,
not just their abstract geometric groups. There is no
change of section origin in (E36.1).

The commutator module of the free class-two pro-three
group on $V$ is $\bigwedge^2V$. The bracket surviving in
(E36.5) is the exterior-square map induced by $V\to T$.
In the basis from (E36.4) it sends $a\wedge b$ to one
generator and sends the other $44$ basic wedges to zero.
Its kernel is therefore exactly $[N,I]/\gamma_3N$,
free of rank $\binom{10}{2}-1=44$; its image has index one.
The orientation identifies the image with
$\bigwedge^2T=\mathbb Z_3(1)$, compatibly with the Weil
pairing. Functoriality proves the asserted arithmetic
action and (E36.2).

We next establish the existence of all the sections in
the statement. Use the rational point $O$ as origin on
the proper elliptic curve. Its full geometric fundamental
group is
$T_\bullet E=\prod_\ell T_\ell E$, and its arithmetic
fundamental group is the corresponding semidirect product
with $G_{\mathbb Q}$. For a given $\alpha$, let
$\beta\in H^1(\mathbb Q,T_3E)$ be its image, and choose
compatible cocycle representatives. There is a specified
full proper-curve section
\[
 g\longmapsto\bigl((\beta(g),0,0,\ldots),g\bigr),
 \tag{E36.7}
\]
where every component at a prime other than three is zero.
The inclusion of the three-primary factor is continuous
and $G_{\mathbb Q}$-equivariant, so (E36.7) is an actual
section of the full arithmetic fundamental group of $E$.

Borne--Emsalem--Stix, Theorem 5.10, states that for a
smooth projective positive-genus curve over $\mathbb Q$
and a boundary which is a union of torsion sub-packets,
every section lifts to its cuspidally central quotient
\cite{E_BorneEmsalemStix}. Apply the theorem to $E$ and
the singleton boundary $\{O\}$. This is a torsion
sub-packet by their Definition 2.3: the only difference
of two points in its support is zero in the Picard group.
All hypotheses are now satisfied. The cuspidal kernel
for a once-punctured elliptic curve is the geometric
commutator subgroup; making it geometrically central
is precisely quotienting by $\gamma_3$. Projecting the
resulting full lift of (E36.7) to its geometric pro-three
quotient gives a section $c_\beta$ of $C\rtimes G_{\mathbb Q}$
above $\beta$.

If this section projects to a representative conjugate
to the chosen $\beta$, conjugate it by a geometric lift
of the relevant element of $T$. Its projection then
equals the chosen representative. Pairing it with the
given abelian section $\alpha$ gives a section of the
fiber product (E36.1). This proves existence for every
$\alpha$, without imposing any ramification condition.

Finally the kernel of $N^{\rm cc}\to V$ is the central
copy of $\mathbb Z_3(1)$ generated by the elliptic
commutator. After choosing one lift of a fixed
representative of $\alpha$, the other lifts differ by
continuous one-cocycles with values in this module.
Central geometric conjugation changes the cocycle by
an ordinary coboundary. There is no additional
identification from geometric elements with nonzero
image in $V$: such an element could fix the prescribed
abelian cocycle only if its image lay in
$H^0(\mathbb Q,V)$, and this group is zero.
Here is a direct verification. In the exact sequence
$0\to T_D\to J_{E,D}\to E\to0$, the rational torsion
of $T_D\simeq\mathbb G_m^2\times
(\operatorname{Res}_{K/\mathbb Q}\mathbb G_m)^3$
is finite, being a product of groups of roots of unity.
The group $E(\mathbb Q)$ has order three by E.24.
Thus $J_{E,D}(\mathbb Q)[3^\infty]$ is finite.
An invariant in $T_3J_{E,D}$ is a compatible sequence
in this finite three-primary group under multiplication
by three. Every such sequence is zero. Hence
$H^0(\mathbb Q,V)=0$ as asserted.

The exact fiber is consequently a torsor under
$H^1(\mathbb Q,\mathbb Z_3(1))$.
Kummer theory identifies this group with
$\varprojlim_m\mathbb Q^\times/\mathbb Q^{\times3^m}$.
Unique factorization in $\mathbb Q^\times$ gives the
isomorphism with (E36.3): the sign is a cube, each
prime valuation is reduced modulo $3^m$, and every
such residue vector has finite support. The inverse
constructs the compatible Kummer classes from precisely
these finite-support residue vectors. This proves the
complete coordinate description of all fibers.

The matrix in (E36.6) has determinant one, and the
exterior-square map in (E36.2) has Smith diagonal one
and kernel rank $44$. The program
\texttt{e\_cuspidally\_central.py} verifies these integral
lattice computations. The existence theorem and the
arithmetic fiber calculation above establish the
infinite section assertions.
\end{proof}
